\SourceLeaf{050}
quae est functio integra indeterminatarum $x,l',l'',l'''$ etc.
\[
\text{per}\quad F(x,l',l'',l'''\text{ etc.})
\]
unde erit identice
\[
\begin{gathered}1+fx\, .\,(\upsilon-Y)+\varphi x\, .\,\frac{\mathrm d(\upsilon-Y)}{\mathrm dx}\\=1+F(x,\lambda',\lambda'',\lambda'''\text{ etc.})\end{gathered}
\]
Habebimus itaque aequationes identicas [1]
\[
\begin{aligned}
&\varphi a\, .\,(a-b)(a-c)(a-d)\, .\, .\, .\, .\\&=1+F(a,\lambda',\lambda'',\lambda'''\text{ etc.})\\
&\varphi b\, .\,(b-a)(b-c)(b-d)\, .\, .\, .\, .\\&=1+F(b,\lambda',\lambda'',\lambda'''\text{ etc.})\\
&\varphi c\, .\,(c-a)(c-b)(c-d)\, .\, .\, .\, .\\&=1+F(c,\lambda',\lambda'',\lambda'''\text{ etc.})\\
&\text{etc.}
\end{aligned}
\]
Supponendo itaque, productum ex omnibus
\[
\begin{gathered}
1+F(a,l',l'',l'''\text{ etc.})\\
1+F(b,l',l'',l'''\text{ etc.})\\
1+F(c,l',l'',l'''\text{ etc.})\\
\text{etc.}
\end{gathered}
\]
quod erit functio integra indeterminatarum $a,b,c$ etc., $l',l'',l'''$ etc. et quidem
functio symmetrica respectu ipsarum $a,b,c$ etc., exhiberi per
\[
\psi(\lambda',\lambda'',\lambda'''\text{ etc.},\ l',l'',l'''\text{ etc.})
\]
e multiplicatione cunctarum aequationum [1] resultabit aequatio identica nova [2]
\[
\pi\varphi a\, .\,\varphi b\, .\,\varphi c\, .\, .\, .\, .=\psi(\lambda',\lambda'',\lambda'''\text{ etc.},\ \lambda',\lambda'',\lambda'''\text{ etc.})
\]
Porro patet, quum productum $\varphi a\, .\,\varphi b\, .\,\varphi c\, .\, .\, .\, .$ indeterminatas $a,b,c$ etc.
symmetrice involvat, inveniri posse functionem integram indeterminatarum
$l',l'',l'''$ etc. talem, quae per substitutiones $l'=\lambda',l''=\lambda'',l'''=\lambda'''$ etc. transeat
in $\varphi a\, .\,\varphi b\, .\,\varphi c\, .\, .\, .\, .$ Sit $t$ illa functio, eritque etiam identice [3]
\[
pt=\psi(l',l'',l'''\text{ etc.},\ l',l'',l''')
\]
quoniam haec aequatio per substitutiones $l'=\lambda',l''=\lambda'',l'''=\lambda'''$ etc. in identicam [2] transit.

Iam ex ipsa definitione functionis $F$ sequitur, identice haberi
\SourceLeaf{051}
\[
F(x,L',L'',L'''\text{ etc.})=0
\]
Hinc etiam identice erit
\[
\begin{gathered}
1+F(a,L',L'',L'''\text{ etc.})=1\\
1+F(b,L',L'',L'''\text{ etc.})=1\\
1+F(c,L',L'',L'''\text{ etc.})=1\\
\text{etc.}
\end{gathered}
\]
et proin erit etiam identice
\[
\psi(\lambda',\lambda'',\lambda'''\text{ etc.},\ L',L'',L'''\text{ etc.})=1
\]
adeoque etiam identice [4]
\[
\psi(l',l'',l'''\text{ etc.},\ L'.\,L'',L'''\text{ etc.})=1
\]
Quamobrem e combinatione aequationum [3] et [4], et substituendo $l'=L'$,
$l''=L'',l'''=L'''$ etc. habebimus [5]
\[
PT=1
\]
si per $T$ denotamus valorem functionis $t$ illis substitutionibus respondentem.
Qui valor quum necessario fiat quantitas finita, $P$ certo nequit esse $=0$. Q. E. D.

\SecondArticle{10}
E praecedentibus iam perspicuum est, quamlibet functionem integram $Y$
unius indeterminatae $x$, cuius determinans sit $=0$, decomponi posse in factores, quorum nullus habeat determinantem 0. Investigato enim divisore communi
altissimo functionum $Y$ et $\frac{\mathrm dY}{\mathrm dx}$, illa iam in duos factores resoluta habebitur. Si
quis horum factorum\HistoricalNote{11} iterum habet determinantem 0, eodem modo in duos factores resolvetur, eodemque pacto continuabimus, donec $Y$ in factores tales tandem
resoluta habeatur, quorum nullus habeat determinantem 0.

Facile porro perspicietur, inter hos factores, in quos $Y$ resolvitur, ad
